\chapter{Yapay Zeka Guvenligi}

\section*{Giris}
Yapay zeka tabanli sistemler hem saldiri yuzeyi hem de savunma kapasitesi olusturur. Bu bolumde AI sistemlerinin korunmasi ve guvenlik operasyonlarinda ureteken AI kullaniminin sinirlari ele alinmaktadir.

\section{Yapay Zeka Sistemlerinin Hedef Alinmasi}
\subsection{Buyuk Dil Modelleri (LLM) Zafiyetleri (Prompt Injection, Veri Zehirlenmesi)}
Prompt injection, data poisoning, model extraction ve yetkisiz arac cagirimi gibi riskler; model etrafinda guvenlik katmanlari ve policy enforcement ile azaltilmalidir.
\begin{itemize}
    \item Girdi/çıkti filtreleme ve policy guardrail
    \item Gizli veri sizintisi icin DLP denetimi
    \item Model ve prompt versiyonlama + denetim izi
\end{itemize}

\section{Siber Guvenlikte Ureteken Yapay Zeka Kullanimi}
Ureteken AI; olay ozetleme, avcilik hipotezi olusturma, playbook taslagi ve kod inceleme gibi alanlarda hiz kazandirir. Ancak insan denetimi ve kanit temelli dogrulama zorunludur.

\Not{AI destekli kararlar "oneri" niteligindedir; olay kapatma ve kritik aksiyonlar icin uzman onayi gereklidir.}
